\chapter[METHODOLOGY]{Methodology}
\label{ch:methodology}

\section{Research Design}
\label{sec:research-design}

This research is quantitative and experimental. Its hypothesis, stated in
Section~\ref{sec:problem-statement}, is that fusing RGB, thermal and LiDAR data
detects corrosion more reliably than any one of these sensors alone. That is a
claim about measurable performance, so it is tested by a controlled comparison:
a set of detectors that differ only in the sensor data they receive, trained and
evaluated under identical conditions, with the metrics and the decision rule
fixed before any model is trained.

The three objectives form a chain. Objective~1 produces the registered trimodal
data without which no fused model can be trained. Objective~2 produces the
detector that consumes those data and the per-region measurements of corroded
area and metal loss. Objective~3 is the comparison itself. It is the only one of
the three whose outcome is open: a finding that a sensor adds nothing to
detection is a valid answer to Research Question~3, and the protocol in
Section~\ref{sec:ablation} is written so that such a finding would be reported.

Two design rules follow from Chapter~\ref{ch:literature-review} and constrain
everything downstream. First, every channel in the fused input comes from a
physical sensor. A channel computed from another channel, such as the texture
map fused with RGB in \citet{Mao2026}, is not a second sensor and is not treated
as one. Second, fusion takes place at the input. Early fusion was the best level
in both reviewed studies that compare fusion levels directly
\citep{Yang2023,Mao2026}, it has been used for corrosion with real sensors
\citep{Lim2021,Maru2025}, and it keeps the ablation clean, because every
configuration uses the same network and only the input channels change. Its
cost is that the sensors must be registered at the pixel level, which is why
registration is an objective in its own right.

\section{Proposed Framework}
\label{sec:framework}

Figure~\ref{fig:framework} shows the proposed framework. The RGB image (three
channels), the thermal image (one channel) and a depth image projected from the
LiDAR point cloud (one channel) are registered to a common view and stacked into
a single five-channel image. One YOLO11 detector \citep{Ultralytics2023} reads
this image and places a bounding box around each corroded region. The LiDAR data
inside each box are then used to measure the corroded area in real units and
the metal loss relative to the surrounding surface.

\begin{figure}[htbp]
\centering
\begin{tikzpicture}[
  scale=1.2, transform shape,
  font=\small,
  box/.style={draw, thick, rounded corners=2pt, align=center, minimum height=1.0cm, inner sep=4pt, fill=white},
  inp/.style={box, minimum width=2.6cm, fill=blue!6},
  core/.style={box, minimum width=2.7cm, fill=blue!60!black, text=white},
  outp/.style={box, minimum width=2.9cm, fill=orange!12},
  arr/.style={-{Stealth[length=2.2mm]}, thick, blue!50!black}
]
  \node[inp] (rgb)  at (0, 1.6)  {RGB image\\\footnotesize 3 channels};
  \node[inp] (th)   at (0, 0)    {Thermal image\\\footnotesize 1 channel};
  \node[inp] (li)   at (0,-1.6)  {LiDAR point cloud\\\footnotesize projected to depth};
  \node[box, minimum width=2.4cm, fill=yellow!15] (reg) at (3.35,0) {Registration\\\footnotesize Objective 1};
  \node[core] (ef)  at (6.45,0)  {Early fusion\\\footnotesize 5-channel image};
  \node[core] (yolo) at (6.45,-2.2) {YOLO11 detector\\\footnotesize pretrained};
  \node[outp] (bb)  at (10.4, 0.7) {Bounding boxes\\\footnotesize corroded regions};
  \node[outp] (am)  at (10.4,-2.2) {Area and metal loss\\\footnotesize per box, from LiDAR};
  \draw[arr] (rgb.east) -- ++(0.35,0) |- (reg.west);
  \draw[arr] (th.east) -- (reg.west);
  \draw[arr] (li.east) -- ++(0.35,0) |- (reg.west);
  \draw[arr] (reg.east) -- (ef.west);
  \draw[arr] (ef.south) -- (yolo.north);
  \draw[arr] (yolo.east) -- ++(0.45,0) |- (bb.west);
  \draw[arr] (bb.south) -- (am.north);
  \draw[arr, dashed] (li.south) -- ++(0,-1.25) -| ([xshift=-6mm]am.south);
\end{tikzpicture}
\caption{Proposed framework. The registered RGB, thermal and LiDAR depth images
are stacked into one five-channel input to a single YOLO11 detector; the LiDAR
data inside each detected box give its corroded area and metal loss.}
\label{fig:framework}
\end{figure}

YOLO11 is chosen for three reasons. YOLO detectors are the most common family in
the reviewed detection studies (Section~\ref{sec:lr-detection}) and in the
box-based corrosion detectors reviewed \citep{Mao2026,Fu2024,Bui2026,Li2024b,Yu2023};
a comparison of YOLOv3, SSD and Faster R-CNN on the same data found YOLOv3 the
fastest with the best results overall \citep{Srivastava2021}, although the
authors note that the best choice depends on the application; and as a
single-stage detector it is fast enough to run on a drone's own computer
\citep{Yu2023}. A YOLO detector has also been trained on RGB, depth and thermal images
fused at the pixel level \citep{Cao2026}. Accepting the five-channel input
requires only that the first convolutional layer be widened. Real-time
deployment on the drone is not an objective of this research.

\section{Project Activities}
\label{sec:project-activities}

The research proceeds through six activities, summarized in
Table~\ref{tab:activities}. Activities 01 and 02 serve Objective~1, Activities 03
and 04 serve Objective~2, and Activities 05 and 06 serve Objective~3.

\begin{table}[htbp]
\centering
\caption{Project activities and the objectives they serve}
\label{tab:activities}
\small
\renewcommand{\arraystretch}{1.2}
\begin{tabularx}{\textwidth}{@{}l>{\RaggedRight\arraybackslash}p{3.6cm}>{\RaggedRight\arraybackslash}Xc@{}}
\toprule
\textbf{No.} & \textbf{Activity} & \textbf{Main output} & \textbf{Objective} \\
\midrule
01 & Data acquisition & RGB, thermal and LiDAR data of corroded metal surfaces, captured from a UAV & 1 \\
02 & Calibration and registration & The three sensors aligned in a common reference frame, with the residual error reported & 1 \\
03 & Preprocessing and annotation & Bounding-box labels and a five-channel input format & 2 \\
04 & Model development & A pretrained YOLO11 detector on the fused input, and per-box area and metal loss from the LiDAR & 2 \\
05 & Ablation protocol & Single-sensor, paired and trimodal arms with a blank-channel control and a fixed decision rule & 3 \\
06 & Evaluation & Detection results against the RGB-only baseline; measurement error against hand measurements & 3 \\
\bottomrule
\end{tabularx}
\end{table}

\subsection{Activity 01: Data Acquisition}
\label{sec:acquisition}

No public dataset pairs registered RGB, thermal and LiDAR data of corroded metal
(Section~\ref{sec:lr-corrosion-detection}). Acquisition is therefore an activity
of this research rather than an input to it, and it is the largest single risk
to the schedule.

The sensors will be carried by an unmanned aerial vehicle (UAV). A UAV can carry
the sensors to structures that are hazardous or slow to reach on foot
\citep{Savino2025,Nooralishahi2022}, and manual inspection of confined
structures is slow: one full inspection of the tanks of an offshore floating
production, storage and offloading unit takes 10 to 15 days \citep{Soares2023}.
The RGB and thermal cameras are already available; the LiDAR is still to be
selected. The three sensors will be mounted rigidly on one frame so that their
relative positions are calibrated once rather than for every capture, the
arrangement used by \citet{Nooralishahi2022} for drone-based inspection of
industrial components.

Data will be collected first from controlled corrosion specimens, whose corroded
areas and metal loss can be measured by hand, and then from field structures if
access is secured. The amount of data needed for an initial study is modest:
\citet{Lim2021} obtained a published RGB--thermal corrosion result from 36
acquired images. Should field access not be secured within the schedule,
evaluation will rely on the specimens alone, and the resulting limit on
generalization will be reported, since models trained on data from a single
source are known to generalize poorly \citep{Nash2022}.

Two acquisition questions are settled in this activity because both affect
whether the method transfers to field inspection.
\begin{itemize}
  \item \textbf{Thermal excitation.} The reviewed thermographic results on
        corrosion were obtained under deliberate thermal excitation: halogen
        lamps at a standoff of 1~m \citep{Lim2021}, a 1000~W quartz-tube heater
        \citep{Elizalde2023} and a freezing spray \citep{Doshvarpassand2022}.
        Whether passive thermal imaging from a UAV reveals the same corrosion is
        not established. Where the specimens allow it, both regimes will be
        recorded on the same specimens so that the difference is measured
        rather than assumed.
  \item \textbf{Timing.} Under active excitation, the thermal measurement is a
        sequence that runs from the excitation until the surface returns to
        ambient temperature \citep{Doshvarpassand2022}, whereas the RGB and LiDAR
        captures are single frames. The acquisition protocol will therefore fix
        which thermal frame is paired with the RGB and LiDAR captures.
\end{itemize}

\subsection{Activity 02: Calibration and Registration}
\label{sec:registration}

The three sensors differ in resolution and field of view, and they respond to
different properties of the scene. A printed checkerboard, the usual target for
camera and LiDAR calibration, has no thermal contrast and is invisible to a
thermal camera. Two routes will therefore be tried. The first is a calibration
target that all three sensors can see, such as the self-heating board of
\citet{Nooralishahi2022}, which uses Peltier modules to produce a sharp pattern
in both thermal and visible images. The second, if the target route fails, is
targetless calibration of the thermal, RGB and LiDAR triple; \citet{Yang2024}
report translation errors of 2.77~cm and 3.86~cm and rotation errors of
0.21\textdegree\ and 0.46\textdegree\ for the LiDAR--RGB and LiDAR--thermal pairs
respectively, and \citet{Tran2024} calibrate the same three sensors without
targets while moving.

After calibration, the LiDAR point cloud is projected onto the RGB image plane to
form the depth image, and the thermal image is warped into the same view.
Registration accuracy will be reported as a result of Objective~1, as the
residual offset in pixels between corresponding features in the registered
images. The preliminary results in Section~\ref{sec:preliminary-results} show
why this matters: a thermal detector lost nearly a third of its accuracy when
scored against boxes drawn on an RGB image that was only approximately aligned
with it.

\subsection{Activity 03: Preprocessing and Annotation}
\label{sec:annotation}

The detection target is a bounding box around each corroded region, with a
single class. Boxes are used instead of pixel masks because building a
pixel-level training set is labor-intensive and time-consuming \citep{Lim2021},
and because mean average precision (mAP) on boxes is the metric reported by the
detection and fusion literature, so the results can be set against published
work. Pixel-level segmentation is outside the scope of this research.

Each channel is normalized separately, because the RGB intensities, the thermal
values and the LiDAR ranges have unrelated scales. Color augmentation that
assumes natural-image statistics, such as hue and saturation shifts, is disabled
in every configuration, because it is not meaningful for a thermal or depth
channel and applying it to some arms and not others would bias the ablation.

No extra labels are needed for the area and metal loss measurements. A box also
covers clean metal, so the corroded pixels inside each box are separated with a
rule-based color method, as in \citet{Alvani2026}, which combines the Lab
$a^*$ channel, a red--green index and an HSV rust map and reached a mean IoU of
0.742 against manually annotated masks, close to the 0.751 of a trained U-Net,
without any pixel labels. A small number of areas will be measured by hand to
state the error of this step.

\subsection{Activity 04: Model Development}
\label{sec:model}

The detector is trained in three steps. It starts from YOLO11 weights pretrained
on the COCO dataset, which teach general visual features and box placement. It
is then fine-tuned on public RGB corrosion images, which teach the appearance of
corrosion; these datasets are too small to train a detector on their own but are
sufficient to adapt a pretrained one. Finally, the first convolutional layer is
widened from three to five channels and the model is trained on the registered
trimodal data: the RGB channels keep their pretrained weights, and the thermal
and depth channels start from new weights.

Transfer learning is used because the trimodal dataset will be small. On steel
pipe pile corrosion, a network trained with pretrained weights converged faster
and was significantly more accurate than the same network trained without them
\citep{Han2024}; a multi-modal bridge deck study notes that training deep
networks from scratch on a limited dataset poses a significant risk of
overfitting \citep{Rashidi2026}; and the LiDAR--image early-fusion study of
section loss used a pretrained network for the same reason \citep{Maru2025}.
Pretraining on RGB images favours the RGB channels, so the comparison is tilted
against fusion and any gain from thermal or LiDAR data is a conservative
estimate. Training from scratch is the fallback if the widened model does not
train well.

The corroded area and metal loss of each detected region are measured from the
LiDAR rather than learned. The range from the LiDAR to the surface converts the
count of corroded pixels in a box into a physical area, and the depth readings
inside the box, compared with the surrounding surface, give the metal loss. The
smallest metal loss that can be measured is set by the precision of the chosen
LiDAR, which will be established when the sensor is selected.

\subsection{Activity 05: Ablation Protocol}
\label{sec:ablation}

The ablation is fixed in full before any corrosion result exists, so that the
comparison to be reported cannot be chosen after the results are seen.

\textbf{Configurations.} The detector is trained on each sensor alone (RGB,
thermal, LiDAR depth), on RGB with thermal, on RGB with LiDAR depth, and on all
three. The single-sensor arms show whether one sensor already suffices; the
paired arms isolate what thermal and LiDAR each add to RGB.

\textbf{Blank-channel control.} One further arm uses the fused input with the
channel under test left empty. A fused input has more channels, and therefore
more input capacity, than a single-sensor one; without this control, a gain due
to the extra capacity would be credited to the sensor.

\textbf{Held constant.} The architecture, pretrained starting weights, optimizer,
learning-rate schedule, number of epochs, early-stopping rule, batch size, image
size, augmentation and data split are identical in every arm. The data are split
by specimen or structure rather than by image, so that frames of the same
corroded region cannot appear in both the training and test sets.

\textbf{Repeats.} Every arm is trained with three random seeds, and results are
reported as the mean over seeds with the observed range.

\textbf{Decision rule.} A sensor earns its place only if the arm containing it
exceeds the best arm without it, on mean mAP@0.5:0.95, by more than the largest
between-seed spread observed in any arm. Comparing the gain with the run-to-run
variation of the experiment itself prevents a difference smaller than the noise
from being reported as an effect. The literature supports both outcomes: an
added sensor improved detection in some studies and not in others
(Section~\ref{sec:lr-fusion}), which is why the question is tested rather than
assumed.

\subsection{Activity 06: Evaluation}
\label{sec:evaluation}

Every configuration is compared with the RGB-only arm, which is the baseline a
fused detector must beat to be worth its cost. Detection is evaluated with
precision, recall, mAP@0.5 and mAP@0.5:0.95. Both mAP thresholds are reported
because thermal and LiDAR data may improve where a box is placed more than
whether a region is found at all: mAP@0.5 counts a box as correct at half
overlap with the label, whereas mAP@0.5:0.95 averages over stricter overlaps and
rewards boxes that fit the corrosion closely. The closest published RGB--thermal
corrosion result \citep{Lim2021} serves as an external reference, with the
differences in structure, detector and thermal excitation stated.

The corroded area and metal loss measured from the LiDAR are compared with hand
measurements on a subset of regions, and their errors are reported in physical
units.

\section{Objectives and Key Milestones}
\label{sec:milestones}

Table~\ref{tab:milestones} links each objective to its activities and to the
milestone that marks its completion.

\begin{table}[htbp]
\centering
\caption{Objectives, activities and key milestones}
\label{tab:milestones}
\small
\renewcommand{\arraystretch}{1.2}
\begin{tabularx}{\textwidth}{@{}>{\RaggedRight\arraybackslash}p{3.4cm}c>{\RaggedRight\arraybackslash}Xl@{}}
\toprule
\textbf{Objective} & \textbf{Activities} & \textbf{Milestone} & \textbf{Target} \\
\midrule
1. Registration and preprocessing & 01, 02 & Registered RGB--thermal--LiDAR dataset ready & Jul 2027 \\
2. Detection and measurement & 03, 04 & Fused detection model finalized & Nov 2027 \\
3. Modality ablation & 05, 06 & Ablation and evaluation complete & Mar 2028 \\
\midrule
-- & -- & Review paper submitted & Feb 2027 \\
-- & -- & Thesis submission and viva & Apr 2028 \\
\bottomrule
\end{tabularx}
\end{table}

\section{Study Plan}
\label{sec:study-plan}

The research runs over six semesters, from May 2026 to April 2028, in six
phases (Figure~\ref{fig:gantt}). The phases overlap by design: the tools are
learned while data acquisition begins, and model development starts before the
dataset is complete. This proposal defence takes place in October 2026, in the
second semester.

\definecolor{gLit}{HTML}{4E79A7}
\definecolor{gProp}{HTML}{3D8C86}
\definecolor{gProg}{HTML}{B07AA1}
\definecolor{gData}{HTML}{59A14F}
\definecolor{gModel}{HTML}{F28E2B}
\definecolor{gThesis}{HTML}{9C755F}
\definecolor{gMile}{HTML}{C0392B}

\begin{landscape}
\begin{figure}[p]
\centering
\begin{ganttchart}[
  hgrid, vgrid,
  x unit=0.66cm,
  y unit title=0.55cm,
  y unit chart=0.62cm,
  title height=1,
  bar height=0.55,
  milestone height=0.55,
  bar/.append style={draw=black!55,thin},
  milestone/.append style={fill=gMile,draw=black!60},
  title label font=\footnotesize,
  bar label font=\small,
  milestone label font=\small\itshape,
  title/.append style={fill=black!8},
  bar label node/.append style={align=left,text width=6.6cm},
  milestone label node/.append style={align=left,text width=6.6cm}
]{1}{24}
\gantttitle{Year 1}{12}
\gantttitle{Year 2}{12} \\
\gantttitle{Semester 1}{4}
\gantttitle{Semester 2}{4}
\gantttitle{Semester 3}{4}
\gantttitle{Semester 4}{4}
\gantttitle{Semester 5}{4}
\gantttitle{Semester 6}{4} \\
\gantttitle{M}{1} \gantttitle{J}{1} \gantttitle{J}{1} \gantttitle{A}{1} \gantttitle{S}{1} \gantttitle{O}{1} \gantttitle{N}{1} \gantttitle{D}{1} \gantttitle{J}{1} \gantttitle{F}{1} \gantttitle{M}{1} \gantttitle{A}{1} \gantttitle{M}{1} \gantttitle{J}{1} \gantttitle{J}{1} \gantttitle{A}{1} \gantttitle{S}{1} \gantttitle{O}{1} \gantttitle{N}{1} \gantttitle{D}{1} \gantttitle{J}{1} \gantttitle{F}{1} \gantttitle{M}{1} \gantttitle{A}{1} \\
\ganttbar[bar/.append style={fill=gLit}]{1. Literature review}{1}{10} \\
\ganttmilestone{Review paper submitted (Feb 2027)}{10} \\
\ganttbar[bar/.append style={fill=gProp}]{2. Proposal}{3}{6} \\
\ganttmilestone{Proposal defence (Oct 2026)}{6} \\
\ganttbar[bar/.append style={fill=gProg}]{3. Tools and coding}{2}{12} \\
\ganttbar[bar/.append style={fill=gData}]{4. Data pipeline}{6}{16} \\
\ganttmilestone{Dataset ready (Jul 2027)}{15} \\
\ganttbar[bar/.append style={fill=gModel}]{5. Model and validation}{13}{23} \\
\ganttmilestone{Model finalized (Nov 2027)}{19} \\
\ganttbar[bar/.append style={fill=gThesis}]{6. Thesis and journal}{17}{24} \\
\ganttmilestone{Thesis submission and viva (Apr 2028)}{24}
\end{ganttchart}
\caption{Study plan, May 2026 to April 2028. Each bar is one phase from start to
finish; diamonds mark the key milestones.}
\label{fig:gantt}
\end{figure}
\end{landscape}

\section{Preliminary Results}
\label{sec:preliminary-results}

Because no public dataset pairs RGB, thermal and LiDAR data of corrosion, the
detection and fusion pipeline was tested on a public trimodal dataset from
another domain. The aim was to show that the pipeline and the ablation protocol
run end to end and produce an interpretable verdict, not to obtain a corrosion
result.

\subsection{Dataset and Setup}

MID-1K \citep{Mota2024} is a person-detection dataset recorded from a Clearpath
Jackal ground robot carrying an RGB camera (Ximea), a thermal camera (FLIR Boson
640) and a 64-beam LiDAR (Ouster OS1-64), with the LiDAR projected onto the RGB
image plane as a depth image. It contains 1,100 frames with one class, person.
The authors' 397 test frames (1,629 labelled people) were kept as the test set,
and the remaining frames were split into 591 training and 112 validation frames,
with the validation frames held out as whole recording sessions.

Every arm used the same detector (YOLO11n, pretrained on COCO), the same split
and the same training recipe: 100 epochs with early stopping after 20 epochs
without improvement, 640-pixel images, a batch size of 16, and hue and
saturation augmentation disabled. Each arm was trained with three seeds. Because
the pretrained detector expects three input channels, the fused arm stacked a
grayscale version of the RGB image with the thermal and depth images, so that
all three sensors fit; the five-channel input of the proposed framework is not
needed for this test. Table~\ref{tab:prelim-arms} lists the six arms.

\begin{table}[htbp]
\centering
\caption{Input configurations in the preliminary test}
\label{tab:prelim-arms}
\small
\renewcommand{\arraystretch}{1.2}
\begin{tabularx}{\textwidth}{@{}l>{\RaggedRight\arraybackslash}p{4.2cm}>{\RaggedRight\arraybackslash}X@{}}
\toprule
\textbf{Arm} & \textbf{Input channels} & \textbf{Scored against} \\
\midrule
RGB & R, G, B & RGB labels \\
Thermal & T, T, T & Thermal labels, drawn on the thermal image \\
Depth & D, D, D & RGB labels \\
Early fusion & Gray, T, D & RGB labels \\
Blank control & Gray, T, empty & RGB labels \\
Thermal on RGB boxes & T, T, T & RGB labels (misalignment test) \\
\bottomrule
\end{tabularx}
\end{table}

\subsection{Results}

Table~\ref{tab:prelim-results} gives the test results.

\begin{table}[htbp]
\centering
\caption{Preliminary test results on MID-1K (397 test frames). Values are the
mean over three seeds, with the range in brackets.}
\label{tab:prelim-results}
\small
\renewcommand{\arraystretch}{1.2}
\begin{tabular}{@{}lll@{}}
\toprule
\textbf{Arm} & \textbf{mAP@0.5} & \textbf{mAP@0.5:0.95} \\
\midrule
\multicolumn{3}{@{}l}{\textit{Scored against RGB labels}} \\
Early fusion (gray + thermal + depth) & 0.596 [0.575--0.608] & 0.324 [0.318--0.328] \\
Blank control (gray + thermal + empty) & 0.598 [0.590--0.605] & 0.324 [0.320--0.332] \\
RGB & 0.572 [0.570--0.575] & 0.311 [0.310--0.313] \\
Thermal on RGB boxes & 0.521 [0.487--0.569] & 0.244 [0.218--0.271] \\
Depth & 0.374 [0.344--0.413] & 0.164 [0.139--0.188] \\
\midrule
\multicolumn{3}{@{}l}{\textit{Scored against thermal labels}} \\
Thermal & 0.756 [0.752--0.763] & 0.467 [0.464--0.471] \\
\bottomrule
\end{tabular}
\end{table}

Early fusion did not satisfy the decision rule. On mAP@0.5:0.95 it reached 0.324
against 0.311 for RGB, a gain of 0.013, while the largest between-seed spread in
any arm was about 0.05. Fusion also scored the same as the blank control, in
which the depth channel was empty, so the LiDAR depth added nothing beyond the
grayscale and thermal channels in this setup. Depth on its own was the weakest
arm.

The thermal images scored 0.756 mAP@0.5 against labels drawn on the thermal
image but 0.521 against the RGB labels, a loss of nearly a third. MID-1K's
thermal camera is only approximately registered to the RGB view. Matching each
RGB label box to the nearest thermal label box gives a median distance of about
5 pixels between box centers, with 90\% of pairs within 13 pixels; this was
measured from the dataset's annotations and includes differences in how the two
label sets were drawn as well as camera misalignment. The fused input is scored
against the RGB labels, so its thermal channel partly points at the wrong place,
which is a likely reason fusion could not exploit the thermal data.

Every arm also scored much higher on the validation set than on the test set,
which reflects the change of scene between the held-out validation sessions and
the authors' test sessions.

\subsection{What the Preliminary Results Establish}

Three lessons carry over to the corrosion study. First, the pipeline runs end to
end: six arms and three seeds were trained and scored with one recipe, and
reruns with the same seed reproduced the same results. Second, the blank-channel
control works: it detected a sensor that added nothing, and the same check will
test whether thermal and LiDAR earn their place on corrosion. Third,
registration has to come first: loose alignment cost thermal nearly a third of
its score on the RGB boxes, which is why calibration and registration form
Objective~1 and Activity~02.

The test does not establish anything about corrosion. The target class is people,
the detector is the smallest YOLO11 model, and the COCO-pretrained weights, which
include a person class, favour the RGB input. Its value is in showing that the
protocol is ready for the trimodal corrosion data.
